\SourceLeaf{196-erchinger}
\Needspace{8\baselineskip}
\pdfbookmark[1]{Erchinger, Siebzehneck. Besprechung vom 19. Dezember 1825}{notice-1825-12-19}
\begin{center}
\rule{.86\linewidth}{.3pt}\par
Göttingische gelehrte Anzeigen.\quad 19. Dezember 1825.\par
\rule{.86\linewidth}{.3pt}
\end{center}
Der Königlichen Gesellschaft ist im Namen von Herrn \textsc{Erchinger} aus Thuningen im Königreich Würtemberg eine kleine Abhandlung vorgelegt worden, die die
\begin{center}\emph{Geometrische Konstruktion des regelmäßigen Siebzehnecks}\end{center}
zum Gegenstand hat. Die allgemeine Theorie der regelmäßigen Vielecke hat bekanntlich durch ihre enge Verbindung mit der höheren Arithmetik eine neue Gestalt und einen erweiterten Umfang erhalten. Ein, wenn auch verhältnismäßig kleiner, Teil davon ist die Theorie der Vielecke, die sich geometrisch konstruieren lassen. Seit der Zeit der Griechen wusste man, dass das Dreieck, das Fünfeck, das Fünfzehneck und alle Vielecke, die daraus durch Verdopplung oder wiederholte Verdopplung der Seitenzahl entstehen, diese Eigenschaft haben; man glaubte und behauptete wohl auch ausdrücklich, dass dies die einzigen seien. Die höhere Arithmetik hat gezeigt, dass dies ein Irrtum war: Indem sie die wahren Quellen der ganz allgemeinen Theorie offenlegte, ergab sich von selbst, dass es außer den genannten Vielecken noch unzählige andere gibt, die sich geometrisch konstruieren lassen und unter denen das Siebzehneck das einfachste ist. Dabei zeigt sich im hellsten Licht die Überlegenheit der Analysis, die das Allgemeinste wie das Besondere mit gleicher Leichtigkeit umfasst, über die Geometrie, die immer beim Besonderen stehen bleiben muss, beim Fortschreiten von einfacheren zu zusammengesetzteren Fällen durch stets wachsende Verwicklungen aufgehalten wird und jenen nächsten Fall nach den bekannten wohl kaum jemals ohne fremde Hilfe erreicht hätte. Dennoch bleibt es wichtig, interessant und wünschenswert, auch rein geometrische Behandlungen fortwährend zu pflegen und der
\SourceLeaf{197}
Geometrie wenigstens einen Teil der neuen Felder zu erschließen, die die Analysis erobert. Dem Rezensenten ist nicht bekannt, dass bisher jemand die Konstruktion des Siebzehnecks öffentlich behandelt hätte, außer Herrn \textsc{Pauker} in den Schriften der Kurländischen Gesellschaft und in seiner Geometrie. Davon verschieden und stärker im rein geometrischen Geist durchgeführt ist die Konstruktion von Herrn \textsc{Erchinger}, die im Folgenden besteht. (Die zugehörige Figur, eine gerade Linie, auf der nacheinander die Punkte $DBGAIFCE$ liegen, kann sich jeder selbst zeichnen.) Man wähle eine beliebige gerade Linie $AB$ und verlängere sie rückwärts und vorwärts nach $C$ und $D$, sodass $AC\times BC=AD\times BD=4AB\times AB$ gilt. Ferner bestimme man die Punkte $E$, $G$ an beiden Seiten der verlängerten Linie $CA$, sodass $AE\times EC=AG\times CG=AB\times AB$, und den Punkt $F$ auf der Seite $A$ der verlängerten Linie $BA$, sodass $AF\times DF=AB\times AB$ gilt. Schließlich teile man $AE$ in $I$, sodass $AI\times EI=AB\times AF$ gilt, wobei $AI$ der kleinere und $EI$ der größere Abschnitt von $AE$ ist. Dann bilde man ein Dreieck, in dem zwei Seiten jeweils $=AB$ und die dritte $=AI$ sind. Beschreibt man um dieses Dreieck einen Kreis, so ist $AI$ die Seite des in diesen Kreis eingeschriebenen regelmäßigen Siebzehnecks.

Prüft man die Richtigkeit dieser Konstruktion durch Vergleich mit der in Art. 354 der \emph{Disquisitiones Arithmeticae} als Beispiel entwickelten Theorie des Siebzehnecks, so erkennt man leicht, dass sie lediglich die geometrische Übersetzung der Gleichungen ist, zu denen die Anwendung der allgemeinen Theorie führt. Tatsächlich sind die Entfernungen der Punkte $C$, $D$, $E$, $F$, $G$, $I$ von $A$ genau die Größen, die dort mit $(8\pdot1)$, $(8\pdot3)$, $(4\pdot1)$, $(4\pdot3)$, $(4\pdot9)$, $(2\pdot1)$ bezeichnet sind, wenn man das positive und negative Zeichen durch die Lage ausdrückt und die Entfernung des Punktes $B$ von $A$ im selben Sinn $=-1$ setzt. Das eigentliche Verdienst der Abhandlung von Herrn \textsc{Erchinger} liegt jedoch weniger in der Aufstellung der Konstruktion selbst, da die Analysis bereits den einfachsten Weg vorgezeichnet hatte, als in der rein geometrischen Begründung ihrer Richtigkeit. Diese ist mit so vorbildlicher, mühevoller Sorgfalt durchgeführt, um alles nicht rein Elementare zu vermeiden, dass sie dem Verfasser zur Ehre gereicht und den Wunsch weckt, sein tatsächlich nicht gewöhnliches mathematisches Talent möge jede Ermutigung finden.
